\section{Where the certificates come from}\label{sec:proof-map}

The probability argument first constructs a threshold $\alpha$ for the
same proper clause law as in Theorem~\ref{thm:main}. It therefore equals
the $k=3$ constant of \emph{A Limiting Satisfiability Threshold for Every
Fixed Clause Size}~\cite[Theorem~1.1]{FixedClauseThreshold}: different
constants would force opposite limits at a density between them. The
comparison identifies the constants. We give the probability argument
locally because the computation needs two additional conclusions:
explicit lower certificates and a positive density of unavoidable
violations above the threshold.

\paragraph{Lower certificates and the sign of pressure.}
For this part of the proof, clauses arrive in a rate-one Poisson process
and their three variable indices are sampled independently, so indices
may repeat within a clause. Let $T_n$ be the first unsatisfiable arrival
time and put $f_n=\E\min\{T_n,20n\}$. Section~\ref{cp:section} proves
\[
 \frac{f_n}{n}\longrightarrow\alpha,
 \qquad \frac{f_n}{n}-2^{67}n^{-1/6}\le\alpha.
\]
Each $f_n$ admits certified rational approximations by finite clause
enumeration and integration. These facts give a computable sequence of
rational lower certificates converging to $\alpha$. Thinning and a
comparison with fixed density slack transfer the threshold from the
auxiliary law to proper clauses.

The estimate behind these conclusions bounds the extra survival time
when one variable may be deleted for free. Its $6/5$ moment is bounded
uniformly even when other deletions are already allowed. The replacement
step follows the residual-set conditioning principle of Friedgut's
half-cube lemma~\cite[Lemma~5.7]{Friedgut1999} and Carenini's sequential
clause replacement~\cite[Theorem~4.3 and Lemma~6.1]{Carenini}. Here
independent variable positions give exact powers for two-literal and
three-literal killing probabilities. A weighted interpolation between
unequal blocks, followed by a weight-balancing argument, turns the moment
bound into a summable error in the normalized mean. Keeping the existing
deletion allowance has a second use: above $\alpha$, every assignment
violates a positive multiple of $n$ clauses with probability tending to
one.

For the auxiliary Poisson formula at time $an$, let $H_n(\sigma)$ count
the clauses violated by an assignment $\sigma$. For a finite penalty
$\beta>0$, define
\[
 Z_n(a,\beta)=\sum_{\sigma\in\{0,1\}^n}e^{-\beta H_n(\sigma)},
 \qquad
 P(a,\beta)=\lim_{n\to\infty}\frac1n\E\log Z_n(a,\beta).
\]
Section~\ref{ctr:section} proves that this pressure limit exists. Below
$\alpha$ it is nonnegative for every $\beta>0$. Above $\alpha$, the
linear violation bound makes it strictly negative for some positive
integer $\beta$.

\paragraph{Finite upper certificates.}
The remaining task is to detect negative pressure using finite data.
Section~\ref{ctr:section} defines an effectively enumerable family of
rational trial descriptions. A description consists of finitely many
rational averaging parameters and a finite table of rational pairs in
$[0,1]^2$. The two coordinates give messages for the two literal signs;
they need not sum to one. Independent uniform coordinates select table
entries through finitely many conditional sampling levels. Different
sites may share initial random data, called the root data, but their
subsequent coordinates are independent. The value $G(a,\beta;Q)$ of a
trial is uniformly computable, and we prove
\[
 P(a,\beta)\le G(a,\beta;Q),\qquad
 P(a,\beta)=\inf_Q G(a,\beta;Q).
\]
A certified negative value is therefore an upper certificate for
$\alpha$, and the infimum identity supplies one at every rational
density above $\alpha$.

The inequality and the infimum identity require different arguments.
For the inequality, the three-literal interpolation leaves the
nonnegative remainder $(x-y)^2(x+2y)$ for $x,y\ge0$. For the reverse
approximation, finite formulas produce nearly optimal trials, but their
later sampling levels can still share random variables between sites.
Removing those variables is the main structural task. Conditional
changes of sampling law and Gaussian tests describe the continuation
law of a fresh site for almost every fixed history of the shared
coordinates. This qualification lets the description survive later
changes of law.

Section~\ref{cs:section} first proves that changing all averaging
parameters has a cost independent of the number of levels. This lets
the tolerances be chosen before the number of structural reductions is
known. Section~\ref{ct:section} then obtains the conditional descriptions
needed by those reductions. Each reduction either reaches the root or
decreases the first remaining parameter by a fixed factor. After
finitely many steps, Section~\ref{cc:section} samples the retained
conditional laws independently at each site, controls the change of the
trial value, and approximates the resulting kernels by finite rational
tables.

Finite hierarchical trials developed in the diluted interpolation and
cavity tradition of Franz and Leone~\cite{FranzLeone2003}, Panchenko
and Talagrand~\cite{PanchenkoTalagrand2004}, and Aizenman, Sims, and
Starr~\cite{AizenmanSimsStarr2003}. The upper-bound and reservoir
calculations below use their interpolation and cavity mechanisms in the
present three-literal model. Panchenko's ultrametricity theorem and the
hierarchical representation of Austin and
Panchenko~\cite{Panchenko2013UM,AustinPanchenko2014} supply the geometry
and representation of the limiting random arrays. Panchenko's pure-state
analysis~\cite{Panchenko2015HE} motivates the identities that also test
the retained site data. The conditional-law records then control
replacement of shared nonroot coordinates by independent site sampling.
This approximation, followed by rational discretization, gives the finite
trial completeness needed for the certificate search.

The proof may select subsequences and conditional laws non-effectively.
Those selections establish that a finite witness exists. The algorithm
searches finite descriptions; it does not receive a selected law or a
convergence modulus. Section~\ref{sec:assembly} closes the argument by
checking the two hypotheses of Lemma~\ref{cert:search}.
